\documentclass{article}
\usepackage[margin=1in]{geometry}
\usepackage{mathtools, amsfonts, amsthm, xcolor, listings}

\title{HW 3}
\author{Sam Ly}

\begin{document}
\maketitle

Final code:

https://github.com/samlyme/Assignments/tree/main

\section*{Chapter 6}

The code for the challenge questions can be found here: 

https://github.com/samlyme/Assignments/tree/cs4080/ch6-challenge

\begin{enumerate}
    \item 
        In C, a block is a statement form that allows you to pack a series of 
        statements where a single one is expected. The comma operator is an 
        analogous syntax for expressions. A comma-separated series of expressions
        can be given where a single expression is expected (except inside a function 
        call’s argument list). At runtime, the comma operator evaluates the left 
        operand and discards the result. Then it evaluates and returns the right
        operand.

        Add support for comma expressions. Give them the same precedence and 
        associativity as in C. Write the grammar, and then implement the 
        necessary parsing code.

        First, to create the grammar, we define a new lowest precedence below
        \verb|equality|. So, the new grammar becomes:
        \begin{verbatim}
expression     → comma;     // CHALLENGE!!
comma         -> equality ("," equality)*;
equality       → comparison ( ( "!=" | "==" ) comparison )* ;
comparison     → term ( ( ">" | ">=" | "<" | "<=" ) term )* ;
term           → factor ( ( "-" | "+" ) factor )* ;
factor         → unary ( ( "/" | "*" ) unary )* ;
unary          → ( "!" | "-" ) unary
               | primary ;
primary        → NUMBER | STRING | "true" | "false" | "nil"
               | "(" expression ")" ;
        \end{verbatim}

        To implement this, we override the current \verb|Expr| method to match 
        for the \verb|Comma| token first.

        \begin{lstlisting}
    private Expr expression() {
        return comma();
    }

    private Expr comma() {
        Expr expr = equality();

        while (match(COMMA)) {
            Token op = previous();
            Expr right = equality();
            expr = new Expr.Binary(expr, op, right);
        }
        return expr;
    }
        \end{lstlisting}


        \item 
            Likewise, add support for the C-style conditional or “ternary” operator ?:. What precedence level is allowed between the ? and :? Is the whole operator left-associative or right-associative?

            The ternary operator has slightly higher precedence than the comma.
            This is because \\ \verb|(a ? b : c, d , e)| should evaluate to 
            \verb|(b, d, e)| or \verb|(c, d, e)|.

            Although the C-style ternary can be nested and is right-associative,
            in my implementation I chose to explicitly enforce the use of parens 
            when nesting ternaries. The implementation also required creating 
            a new Expression type since the ternary expression is actually a 
            trinary expression.

            \begin{lstlisting}
    private Expr ternary() {
        Expr expr = equality();

        if (match(QUESTION)) {
            Token leftOp = previous();
            Expr center = equality();

            consume(COLON, "Expect ':' in ternary.");
            Token rightOp = previous();
            Expr right = equality();

            expr = new Expr.Trinary(expr, leftOp, center, rightOp, right);
        }

        return expr;
    }
            \end{lstlisting}

        \item 
            Add error productions to handle each binary operator appearing without a left-hand operand. In other words, detect a binary operator appearing at the beginning of an expression. Report that as an error, but also parse and discard a right-hand operand with the appropriate precedence.

            The grammar for error productions now looks like this:
            \begin{verbatim}
expression     → comma;     // CHALLENGE!!
comma         -> ternary ("," ternary)*;
ternary       -> ternary "?" ternary ":" ternary;
               | equality;
equality       → comparison? ( ( "!=" | "==" ) comparison )* ;
comparison     → term? ( ( ">" | ">=" | "<" | "<=" ) term )* ;
term           → factor? ( ( "-" | "+" ) factor )* ;
factor         → unary? ( ( "/" | "*" ) unary )* ;
unary          → ( "!" | "-" ) unary
               | primary ;
primary        → NUMBER | STRING | "true" | "false" | "nil"
               | "(" expression ")" ;
            \end{verbatim}

            However, when implementing, we defer the check all the down to the 
            \verb|primary| method because it allows us to handle it as an error.

            \begin{verbatim}
    private Expr primary() {
        if (match(FALSE)) return new Expr.Literal(false);
        if (match(TRUE)) return new Expr.Literal(true);
        if (match(NIL)) return new Expr.Literal(null);

        if (match(NUMBER, STRING)) {
            return new Expr.Literal(previous().literal);
        }

        if (match(LEFT_PAREN)) {
            Expr expr = expression();
            consume(RIGHT_PAREN, "Expect ')' after expression.");
            return new Expr.Grouping(expr);
        }

        if (match(PLUS, MINUS, STAR, SLASH, 
                  BANG_EQUAL, EQUAL_EQUAL, GREATER, 
                  GREATER_EQUAL, LESS, LESS_EQUAL)) {
            error(previous(), "Missing left operand");
            term();
            return  null;
        }

        throw error(peek(), "Expect expression.");
    }
            \end{verbatim}

\end{enumerate}

\section*{Chapter 7}
The code for the challenge questions can be found here: 

https://github.com/samlyme/Assignments/tree/cs4080/ch7/challenge
\begin{enumerate}
    \item {
        Allowing comparisons on types other than numbers could be useful. The operators might have a reasonable interpretation for strings. Even comparisons among mixed types, like 3 < "pancake" could be handy to enable things like ordered collections of heterogeneous types. Or it could simply lead to bugs and confusion.

        Would you extend Lox to support comparing other types? If so, which pairs of types do you allow and how do you define their ordering? Justify your choices and compare them to other languages.

        I personally would not extend Lox to support comparing other types for 
        two reasons: 1) ambiguity and unexpected behavior. Lox is already 
        dynamically typed, so incorporating this could lead to cursed situations 
        where you need to reason about the hypotheical behavior of an operator 
        based on the hypotheical type of its operands. 2) implementation difficulty.
        There is no good way to perform dynamic dispatch at this level. Sure, we 
        can override each `internal' type to have a custom comparing method, but that
        method would then need to somehow infer the `absolute value' of arbitrary 
        objects. A good example of this is Java itself! Java explicitly does not 
        allow for comparing different types.
    }
    \item {
        Many languages define + such that if either operand is a string, the other is converted to a string and the results are then concatenated. For example, "scone" + 4 would yield scone4. Extend the code in visitBinaryExpr() to support that.

        \begin{lstlisting}
            case PLUS:
                if (left instanceof Double && right instanceof Double) {
                    return (double)left + (double)right;
                }
                if (left instanceof String || right instanceof String) {
                    return stringify(left) + stringify(right);
                }
                throw new RuntimeError(expr.operator,
                        "Operands must be two numbers or two strings.");
        \end{lstlisting}
    }

    \item {
        What happens right now if you divide a number by zero? What do you think should happen? Justify your choice. How do other languages you know handle division by zero, and why do they make the choices they do?

        Change the implementation in visitBinaryExpr() to detect and report a runtime error for this case.

        Now, if we divide by zero, we get `Infinity'. This is a natural consequence 
        of Java's implementation of the Double class. I think this actually makes 
        sense! It gives a good way for the program to gracefully handle edge cases
        that is also obvious for users to debug. Some other languages either return 
        a special `NaN' value or throw an exception. These choices also make sense.
        By returning `NaN', the downstream mathematical operations can be halted,
        and also gives the user a good idea of the source of the bug. By throwing
        an exception, the user is forced to handle the edge case. 

        \begin{lstlisting}
            case SLASH:
                checkNumberOperands(expr.operator, left, right);
                if (right.equals(0.0)) 
                    throw new RuntimeError(
                        expr.operator, "Can't divide by zero.");
                return (double)left / (double)right;
        \end{lstlisting}

    }
\end{enumerate}
\end{document}
